% theory_endogenous_premium.tex — Endogenous Premium Theory
% Core theoretical contribution: the NAV premium as a perpetual accumulation option

\section{Endogenous Premium Theory}\label{sec:endogenous_premium}

The methodology of Section~\ref{sec:methodology} treats the NAV premium
$\pi_t$ as an exogenous Ornstein--Uhlenbeck process with fixed long-run mean
$\theta$.  This is descriptively adequate but theoretically unsatisfying: it
does not explain \emph{why} a rational market would price a BTC treasury
vehicle above (or below) its net asset value.  This section develops an
endogenous theory of the premium.  The central insight is that the NAV premium
reflects the market value of a \emph{perpetual accumulation option}---the
option to issue equity at a premium and convert proceeds into BTC, thereby
increasing BTC-per-share for existing holders.

We proceed in four parts.  Section~\ref{subsec:endo_system} specifies the
coupled stochastic system that replaces the exogenous OU.
Section~\ref{subsec:fair_premium} derives the equilibrium (fair) premium
$\pi^*$ as the present value of future accretive issuance.
Section~\ref{subsec:reflexivity} characterizes the reflexive feedback loop and
establishes conditions for multiple equilibria.
Section~\ref{subsec:mispricing} defines the mispricing metric and analyzes its
dynamics.

%% ====================================================================
\subsection{The Endogenous Premium System}\label{subsec:endo_system}

\subsubsection{Motivation}

Under the exogenous specification $d\pi_t = \kappa(\theta - \pi_t)\,dt +
\sigma_\pi\,dW_t^\pi$, the long-run mean $\theta$ is a free parameter
estimated from data.  But $\theta$ should itself depend on the firm's
fundamentals: its BTC holdings, leverage, volatility environment, and
capacity to issue accretively.  We now endogenize $\theta$ by replacing it
with a state-dependent equilibrium premium $\pi^*(\Phi_t)$.

\subsubsection{State Vector and Assumptions}

\begin{assumption}[Endogenous Premium Environment]\label{ass:endo_env}
We work under the following conditions:
\begin{enumerate}
  \item The BTC spot price $S_t$ follows geometric Brownian motion under
    $\mathbb{P}$ with drift $\mu_S$ and volatility $\sigma_S$, augmented by
    a market impact term from the firm's BTC purchases.
  \item The firm's BTC holdings $H_t$, debt $D_t$, preferred liabilities
    $P_t^{\mathrm{pfd}}$, and share count $N_t^{\mathrm{sh}}$ evolve as in
    Sections~\ref{sec:methodology}--\ref{sec:model_extended}.
  \item The firm issues equity (ATM) whenever $\pi_t$ exceeds the accretion
    threshold $\pi^*_{\mathrm{ATM}} = \log(\mathrm{ELE}_t)$ from
    Proposition~\ref{prop:accretion}, at a rate proportional to
    $(\pi_t - \pi^*_{\mathrm{ATM}})^+$.
  \item Market participants are risk-neutral (for valuation of the
    accumulation option) and observe the full state vector.
  \item The risk-free rate $r > 0$ is constant.
\end{enumerate}
\end{assumption}

\begin{definition}[State Vector]\label{def:state_vector}
The endogenous premium model is driven by the state vector
\begin{equation}\label{eq:state_vector}
  \Phi_t = \bigl(S_t,\, H_t,\, D_t,\, P_t^{\mathrm{pfd}},\,
  N_t^{\mathrm{sh}},\, \sigma_S,\, \kappa,\, r\bigr),
\end{equation}
from which we derive the ancillary quantities $A_t = H_t S_t$,
$L_t = D_t + P_t^{\mathrm{pfd}}$, $\mathrm{NAV}_t = (A_t - L_t)^+$,
$B_t = H_t / N_t^{\mathrm{sh}}$, and
$\mathrm{ELE}_t = A_t / (A_t - L_t)$.
\end{definition}

\subsubsection{The Coupled SDE System}

\begin{definition}[Endogenous Premium Dynamics]\label{def:endo_sdes}
The endogenous premium model consists of the following coupled system of
stochastic differential equations:
\begin{align}
  \frac{dS_t}{S_t} &= \mu_S\, dt + \sigma_S\, dW_t^S
    + \eta \cdot \frac{dH_t}{H_t},
    \label{eq:dS_endo} \\[6pt]
  d\pi_t &= \kappa\bigl(\pi^*(\Phi_t) - \pi_t\bigr)\, dt
    + \sigma_\pi\, dW_t^\pi,
    \label{eq:dpi_endo} \\[6pt]
  dH_t &= \alpha\, \bigl(\pi_t - \log(\mathrm{ELE}_t)\bigr)^+\, dt
    + dH_t^{\mathrm{jump}},
    \label{eq:dH_endo}
\end{align}
where:
\begin{itemize}
  \item $\eta \ge 0$ is the \emph{market impact coefficient}: each unit of
    BTC purchased by the firm pushes the BTC price by a factor $\eta$.
  \item $\pi^*(\Phi_t)$ is the \emph{equilibrium (fair) premium}, derived
    in Section~\ref{subsec:fair_premium}, replacing the exogenous $\theta$.
  \item $\kappa > 0$ is the mean-reversion speed toward the fair premium.
  \item The accumulation rate in \eqref{eq:dH_endo} is driven by the
    \emph{excess premium} $(\pi_t - \log(\mathrm{ELE}_t))^+$ rather than
    $\pi_t^+$, reflecting that only accretive issuance
    (Proposition~\ref{prop:accretion}) is undertaken.
  \item $\mathrm{corr}(dW_t^S, dW_t^\pi) = \rho$.
\end{itemize}
\end{definition}

\begin{remark}[Comparison with Exogenous Model]\label{rem:comparison}
The exogenous model of Section~\ref{sec:methodology} is recovered by setting
$\eta = 0$ (no market impact), $\pi^*(\Phi_t) = \theta$ (constant long-run
mean), and replacing $(\pi_t - \log(\mathrm{ELE}_t))^+$ with $\pi_t^+$ in
the holding dynamics.  The endogenous model thus nests the exogenous
specification as a special case.
\end{remark}

\begin{remark}[Reflexive Structure]\label{rem:reflexive}
The system \eqref{eq:dS_endo}--\eqref{eq:dH_endo} exhibits Soros-type
reflexivity.  A high premium $\pi_t$ triggers accumulation ($dH_t > 0$),
which via market impact raises $S_t$, which raises $\mathrm{NAV}_t$, which
(if $\pi^*$ increases with fundamentals) raises the fair premium, which
sustains the high observed premium.  The feedback is self-reinforcing in
the upward direction and self-reinforcing in the downward direction, creating
the possibility of multiple equilibria (Section~\ref{subsec:reflexivity}).
\end{remark}

%% ====================================================================
\subsection{Fair Premium Derivation}\label{subsec:fair_premium}

We now derive $\pi^*(\Phi_t)$---the premium that a rational market would
assign to a BTC treasury vehicle given its current fundamentals and
accumulation capacity.

\subsubsection{The Accumulation Option}

\begin{definition}[Accumulation Option]\label{def:accum_option}
The \emph{accumulation option} is the perpetual right to issue equity at a
premium to NAV and convert proceeds into BTC, thereby increasing BTC-per-share.
Its value to existing shareholders at time $t$ is:
\begin{equation}\label{eq:V_accretion}
  V_{\mathrm{acc}}(t) = \mathbb{E}_t\!\left[
    \int_t^\infty e^{-r(u-t)}\,
    \underbrace{\bigl(\pi_u - \log(\mathrm{ELE}_u)\bigr)^+}_{\text{excess premium}}
    \cdot\, \underbrace{q(u)\, S_u}_{\text{BTC value acquired}}
    \cdot\, \underbrace{\frac{N_t^{\mathrm{sh}}}{N_u^{\mathrm{sh}}}
    }_{\text{dilution adjustment}}\,
    du
  \right],
\end{equation}
where $q(u)$ is the rate of BTC acquisition at time $u$ (in BTC per unit time)
and the dilution adjustment accounts for the fact that future issuance dilutes
existing shareholders.
\end{definition}

\begin{remark}
The integral \eqref{eq:V_accretion} has the structure of a perpetual
exchange option: at each instant, the firm can ``exchange'' newly issued
shares (at market price) for BTC (at spot), and the option is in-the-money
whenever the equity premium exceeds the leverage threshold.  The perpetual
nature---there is no expiry on the ATM program---makes this more valuable
than a finite-horizon option.
\end{remark}

\subsubsection{Simplified Closed-Form Derivation}

To obtain a tractable expression for $\pi^*$, we work under a set of
simplifying assumptions.

\begin{assumption}[Stationary Accumulation]\label{ass:stationary}
\begin{enumerate}
  \item The premium $\pi_t$ follows a stationary OU process around $\pi^*$
    with speed $\kappa$ and volatility $\sigma_\pi$.
  \item BTC holdings grow at a constant proportional rate $g_H$ (the
    model-implied IBGR from Section~\ref{sec:methodology}).
  \item Dilution offsets a fraction $\delta$ of gross BTC growth, so the
    net per-share BTC growth rate is $g_B = g_H - \delta$.
  \item The leverage ratio $\ell_t = L_t / A_t$ is approximately constant
    at $\bar{\ell}$.
  \item Market impact is small: $\eta \ll 1$.
\end{enumerate}
\end{assumption}

\begin{proposition}[Accretive Issuance Value]\label{prop:accretion_value}
Under Assumptions~\ref{ass:endo_env} and~\ref{ass:stationary}, the
per-share value of the accumulation option is
\begin{equation}\label{eq:V_acc_closed}
  V_{\mathrm{acc}} = \mathrm{NAV}_t^{\mathrm{ps}} \cdot
  \frac{g_B}{r - g_B} \cdot
  \Psi(\pi^*, \kappa, \sigma_\pi, \bar{\ell}),
\end{equation}
where $\mathrm{NAV}_t^{\mathrm{ps}} = \mathrm{NAV}_t / N_t^{\mathrm{sh}}$
is per-share NAV, $g_B$ is the net per-share BTC growth rate, and
\begin{equation}\label{eq:Psi}
  \Psi(\pi^*, \kappa, \sigma_\pi, \bar\ell) =
  \Phi_{\mathrm{N}}\!\left(
    \frac{\pi^* - \log(1/(1-\bar\ell))}{\sigma_\pi / \sqrt{2\kappa}}
  \right)
  \cdot \left(1 + \frac{\sigma_\pi^2}{4\kappa\, r}\right)
\end{equation}
is the \emph{accumulation efficiency factor}.  Here
$\Phi_{\mathrm{N}}(\cdot)$ is the standard normal CDF.
\end{proposition}

\begin{proof}
Under Assumption~\ref{ass:stationary}, the premium fluctuates around $\pi^*$
with stationary distribution
$\pi \sim \mathcal{N}(\pi^*,\, \sigma_\pi^2 / (2\kappa))$.  The accumulation
option is exercised whenever $\pi_t > \log(\mathrm{ELE}_t) =
\log(1/(1-\bar\ell))$.  The probability that the option is in-the-money at
any given time is
\begin{equation}\label{eq:prob_itm}
  p_{\mathrm{ITM}} = \mathbb{P}\!\left(
    \pi_t > \log\!\left(\frac{1}{1-\bar\ell}\right)
  \right)
  = \Phi_{\mathrm{N}}\!\left(
    \frac{\pi^* - \log(1/(1-\bar\ell))}{\sigma_\pi / \sqrt{2\kappa}}
  \right).
\end{equation}

When in-the-money, the expected excess premium is proportional to
$\sigma_\pi / \sqrt{2\kappa}$ (the standard deviation of the stationary
distribution), contributing the correction factor $(1 + \sigma_\pi^2 /
(4\kappa r))$ after integrating the discounted expected payoff.

The per-share value of BTC accretion is a perpetuity growing at rate $g_B$,
discounted at rate $r$:
\[
  \mathrm{NAV}_t^{\mathrm{ps}} \cdot \frac{g_B}{r - g_B}
\]
whenever $g_B < r$ (transversality condition).  Multiplying by the fraction
of time the option is exercisable and the volatility correction yields
\eqref{eq:V_acc_closed}.
\end{proof}

\subsubsection{The Fair Premium}

\begin{theorem}[Fair Premium]\label{thm:fair_premium}
Under Assumptions~\ref{ass:endo_env} and~\ref{ass:stationary}, the
equilibrium log-premium $\pi^*$ satisfies the fixed-point equation
\begin{equation}\label{eq:pi_star_fp}
  \boxed{
    \pi^* = \log\!\left(1 + \frac{V_{\mathrm{acc}}}{\mathrm{NAV}_t^{\mathrm{ps}}}\right)
    = \log\!\left(1 + \frac{g_B}{r - g_B}\,
    \Psi(\pi^*, \kappa, \sigma_\pi, \bar\ell)\right).
  }
\end{equation}
\end{theorem}

\begin{proof}
The equity market capitalization per share equals NAV per share plus the
value of the accumulation option:
\begin{equation}\label{eq:equity_decomp_endo}
  \frac{E_t}{N_t^{\mathrm{sh}}}
  = \mathrm{NAV}_t^{\mathrm{ps}} + V_{\mathrm{acc}}.
\end{equation}
The log-premium is $\pi_t = \log(E_t / \mathrm{NAV}_t)$, so in equilibrium:
\begin{align}
  e^{\pi^*} &= \frac{E_t / N_t^{\mathrm{sh}}}{\mathrm{NAV}_t / N_t^{\mathrm{sh}}}
  = \frac{\mathrm{NAV}_t^{\mathrm{ps}} + V_{\mathrm{acc}}}
         {\mathrm{NAV}_t^{\mathrm{ps}}}
  = 1 + \frac{V_{\mathrm{acc}}}{\mathrm{NAV}_t^{\mathrm{ps}}}.
  \label{eq:pi_star_derivation}
\end{align}
Substituting \eqref{eq:V_acc_closed} and taking logarithms yields
\eqref{eq:pi_star_fp}.

This is a fixed-point equation because $\Psi$ depends on $\pi^*$ through
the normal CDF term in \eqref{eq:Psi}.  Existence and uniqueness are
established in Proposition~\ref{prop:fp_existence}.
\end{proof}

\begin{proposition}[Existence and Uniqueness of Fair Premium]%
\label{prop:fp_existence}
Define the map
$F(\pi) = \log\!\bigl(1 + (g_B/(r-g_B))\,\Psi(\pi, \kappa, \sigma_\pi,
\bar\ell)\bigr)$.  Under the condition $g_B < r$:
\begin{enumerate}
  \item $F$ maps $\mathbb{R}$ into $[0, \bar\pi]$ where
    $\bar\pi = \log(1 + g_B/(r-g_B)) < \infty$.
  \item $F$ is increasing and Lipschitz with constant
    $\mathrm{Lip}(F) < 1$ whenever
    \begin{equation}\label{eq:contraction_cond}
      \frac{g_B}{r - g_B} \cdot
      \frac{\phi_{\mathrm{N}}(\cdot)}{\sigma_\pi / \sqrt{2\kappa}}
      \cdot \frac{1}{1 + g_B \Psi / (r - g_B)} < 1,
    \end{equation}
    where $\phi_{\mathrm{N}}$ is the standard normal PDF.
  \item Under \eqref{eq:contraction_cond}, $F$ is a contraction mapping on
    $[0, \bar\pi]$, so by the Banach fixed-point theorem there exists a
    unique $\pi^* \in [0, \bar\pi]$ satisfying $F(\pi^*) = \pi^*$.
\end{enumerate}
\end{proposition}

\begin{proof}
\emph{(i)} Since $\Psi \in [0, 1 + \sigma_\pi^2/(4\kappa r)]$ and
$g_B/(r-g_B) > 0$, we have $F(\pi) \in [0, \bar\pi]$ for all $\pi$.

\emph{(ii)} The derivative of $F$ is
\[
  F'(\pi) = \frac{g_B/(r-g_B)}{1 + g_B\,\Psi(\pi)/(r-g_B)}
  \cdot \frac{\partial \Psi}{\partial \pi}.
\]
Since $\partial\Psi/\partial\pi = \phi_{\mathrm{N}}(\cdot) /
(\sigma_\pi/\sqrt{2\kappa}) \cdot (1 + \sigma_\pi^2/(4\kappa r))$, the
Lipschitz constant is bounded by the supremum of $|F'|$, which is finite.
Condition \eqref{eq:contraction_cond} ensures $\sup|F'| < 1$.

\emph{(iii)} A contraction mapping on a complete metric space (the closed
interval $[0,\bar\pi] \subset \mathbb{R}$) has a unique fixed point by
the Banach theorem.
\end{proof}

\subsubsection{Comparative Statics}

\begin{proposition}[Comparative Statics of the Fair Premium]%
\label{prop:comp_statics}
The equilibrium premium $\pi^*$ satisfies the following comparative statics:
\begin{align}
  \frac{\partial \pi^*}{\partial g_B} &> 0
    &\quad& \text{(higher per-share BTC growth $\Rightarrow$ higher premium)},
    \label{eq:cs_gB} \\[4pt]
  \frac{\partial \pi^*}{\partial \sigma_S} &> 0
    &\quad& \text{(higher BTC volatility $\Rightarrow$ higher premium)},
    \label{eq:cs_sigma} \\[4pt]
  \frac{\partial \pi^*}{\partial \kappa} &< 0
    &\quad& \text{(faster mean-reversion $\Rightarrow$ lower premium)},
    \label{eq:cs_kappa} \\[4pt]
  \frac{\partial \pi^*}{\partial \bar\ell} &\lessgtr 0
    &\quad& \text{(leverage: non-monotone)},
    \label{eq:cs_ell} \\[4pt]
  \frac{\partial \pi^*}{\partial r} &< 0
    &\quad& \text{(higher discount rate $\Rightarrow$ lower premium)}.
    \label{eq:cs_r}
\end{align}
\end{proposition}

\begin{proof}
We apply the implicit function theorem to $G(\pi^*, \theta) =
\pi^* - F(\pi^*; \theta) = 0$, where $\theta$ denotes a generic parameter.
Since $\partial G / \partial \pi^* = 1 - F'(\pi^*) > 0$ by the contraction
condition, we have
\begin{equation}\label{eq:ift}
  \frac{\partial \pi^*}{\partial \theta}
  = \frac{\partial F / \partial \theta}{1 - F'(\pi^*)}.
\end{equation}
The sign is determined by $\partial F / \partial \theta$:

\emph{BTC growth rate $g_B$:} $\partial F / \partial g_B > 0$ since
higher growth increases both $g_B/(r-g_B)$ and the present value of future
accretion.

\emph{BTC volatility $\sigma_S$:} Higher $\sigma_S$ increases BTC price
dispersion, which increases the expected future NAV growth (via the convexity
of the equity claim on BTC) and raises the value of the accumulation option.
This effect operates through $g_B$, which is increasing in $\sigma_S$ via
the flywheel's option-like payoff.

\emph{Mean-reversion speed $\kappa$:} Higher $\kappa$ reduces the
stationary variance $\sigma_\pi^2/(2\kappa)$ of the premium, making extreme
premium excursions (which drive the most accretive issuance) less likely.
This lowers $\Psi$ and hence $\pi^*$.

\emph{Leverage $\bar\ell$:} The effect is non-monotone.  Higher leverage
raises the accretion threshold $\log(1/(1-\bar\ell))$, reducing the
probability that issuance is accretive (lowering $p_{\mathrm{ITM}}$).
However, higher leverage also amplifies BTC returns via $\mathrm{ELE}_t$,
potentially increasing $g_B$.  At low leverage, the amplification effect
dominates; at high leverage, the threshold effect dominates, creating an
interior optimum.

\emph{Discount rate $r$:} Higher $r$ reduces the present value of the
perpetuity $g_B/(r-g_B)$, directly lowering $\pi^*$.
\end{proof}

\begin{remark}[Volatility as Value Driver]\label{rem:vol_value}
The comparative static $\partial\pi^*/\partial\sigma_S > 0$ is one of the
model's most distinctive predictions: unlike a standard equity, where higher
volatility increases risk and reduces value (all else equal), a BTC treasury
vehicle \emph{benefits} from higher BTC volatility because it makes the
accumulation option more valuable.  This is analogous to the vega of a call
option.  Empirically, MSTR's premium has tended to expand during periods of
elevated BTC volatility, consistent with this prediction.
\end{remark}

\subsubsection{Conditions for Premium vs.\ Discount}

\begin{proposition}[Premium--Discount Boundary]\label{prop:premium_discount}
The fair premium satisfies $\pi^* > 0$ (premium to NAV) if and only if the
accumulation option has positive value, which requires:
\begin{equation}\label{eq:premium_condition}
  g_B > 0
  \quad\text{and}\quad
  p_{\mathrm{ITM}} > 0
  \quad\text{and}\quad
  g_B < r.
\end{equation}
The fair premium satisfies $\pi^* < 0$ (discount to NAV) when either:
\begin{enumerate}
  \item $g_B < 0$: dilution exceeds gross BTC accumulation, destroying
    per-share value; or
  \item The firm faces a perpetual cash drain (preferred dividends, debt
    coupons) that exceeds expected asset appreciation, so the ongoing
    liabilities reduce NAV faster than the flywheel builds it.
\end{enumerate}
\end{proposition}

\begin{proof}
From \eqref{eq:pi_star_fp}, $\pi^* > 0$ iff
$V_{\mathrm{acc}} / \mathrm{NAV}_t^{\mathrm{ps}} > 0$.  By
\eqref{eq:V_acc_closed}, $V_{\mathrm{acc}} > 0$ requires $g_B > 0$
(positive net accretion), $\Psi > 0$ (positive probability of accretive
issuance), and $r > g_B$ (finite present value).  All three conditions
in \eqref{eq:premium_condition} are necessary and sufficient.

For the discount case, when the firm faces perpetual liabilities with
present value exceeding the accumulation option, we extend
\eqref{eq:equity_decomp_endo} to
\[
  \frac{E_t}{N_t^{\mathrm{sh}}}
  = \mathrm{NAV}_t^{\mathrm{ps}} + V_{\mathrm{acc}}
  - \mathrm{PV}(\text{net liability drag}),
\]
where the net liability drag is the present value of preferred dividends
and coupons net of expected asset appreciation from the flywheel.
When $V_{\mathrm{acc}} < \mathrm{PV}(\text{drag})$, the premium turns
negative.
\end{proof}

\begin{remark}[Interpreting NAV Discounts]
A persistent NAV discount ($\pi^* < 0$) signals that the market views the
firm's capital structure as value-destroying: the cost of servicing
liabilities exceeds the benefit of the accumulation option.  This can occur
when (i) the premium is too low for accretive ATM issuance, (ii) convertible
debt conversion dilutes without commensurate BTC growth, or (iii) preferred
dividend obligations ($\Phi_t$) consume expected returns.  The discount is
\emph{rational} in this framework, not a market inefficiency.
\end{remark}

%% ====================================================================
\subsection{Reflexivity and Multiple Equilibria}\label{subsec:reflexivity}

The coupled system \eqref{eq:dS_endo}--\eqref{eq:dH_endo} exhibits
self-reinforcing dynamics that can generate multiple equilibria.  We now
formalize this mechanism.

\subsubsection{The Feedback Map}

\begin{definition}[Feedback Map]\label{def:feedback_map}
Define the \emph{premium feedback map}
$\mathcal{F}: \pi \mapsto \pi'$ as the composition of:
\begin{enumerate}
  \item \textbf{Accumulation:} Premium $\pi$ induces BTC acquisition at
    rate $q(\pi) = \alpha(\pi - \log(\mathrm{ELE}))^+$.
  \item \textbf{Market impact:} BTC purchases push the BTC price by
    $\Delta S / S = \eta \cdot q(\pi) / H$.
  \item \textbf{NAV update:} Higher $S$ raises $\mathrm{NAV}$ and
    potentially changes the fair premium.
  \item \textbf{Premium update:} The market re-prices the accumulation
    option given the updated state, yielding $\pi' = \pi^*(\Phi')$.
\end{enumerate}
An equilibrium is a fixed point of $\mathcal{F}$: $\pi = \mathcal{F}(\pi)$.
\end{definition}

\begin{theorem}[Multiple Equilibria]\label{thm:multiple_eq}
Under Assumptions~\ref{ass:endo_env}--\ref{ass:stationary}, if the market
impact parameter $\eta$ is sufficiently large and the accumulation rate
$\alpha$ is sufficiently responsive to the premium, there exist (generically)
three equilibria:
\begin{enumerate}
  \item \textbf{High-premium equilibrium} $\pi_H^* > 0$: The premium is
    high, the firm issues accretively, BTC accumulation supports the BTC
    price, and the premium is self-sustaining.  This equilibrium is
    \emph{locally stable}.
  \item \textbf{Low-premium (discount) equilibrium} $\pi_L^* \le 0$: The
    premium is at or below zero, the firm cannot issue accretively, there is
    no BTC accumulation, no market impact support, and the discount is
    self-reinforcing.  This equilibrium is \emph{locally stable}.
  \item \textbf{Unstable equilibrium} $\pi_U^* \in (\pi_L^*, \pi_H^*)$:
    An intermediate premium that is a saddle point---any perturbation pushes
    the system toward either $\pi_H^*$ or $\pi_L^*$.
\end{enumerate}
\end{theorem}

\begin{proof}
Consider the net premium adjustment function
$G(\pi) = \mathcal{F}(\pi) - \pi$, which measures the difference between
the implied fair premium and the current premium.  An equilibrium satisfies
$G(\pi) = 0$.

\emph{Step 1: Shape of $G$.}
At $\pi \to -\infty$, no issuance occurs, so $q(\pi) = 0$, there is no
market impact, and the fair premium is determined solely by the (negative)
liability drag: $\mathcal{F}(\pi) = \pi_0^{\mathrm{base}} > -\infty$, a
finite value.  Thus $G(\pi) > 0$ for sufficiently negative $\pi$.

At $\pi \to +\infty$, issuance is massive, but per-share dilution grows
proportionally, and the marginal value of additional BTC decreases (each
additional BTC is purchased at an increasingly high dollar cost due to
market impact).  Thus $\mathcal{F}(\pi)$ grows slower than $\pi$, and
$G(\pi) < 0$ for sufficiently large $\pi$.

\emph{Step 2: Non-monotonicity.}
For intermediate $\pi$, the feedback loop creates a region where
$\mathcal{F}(\pi)$ exceeds $\pi$: the market impact of accumulation raises
the BTC price, which raises NAV, which raises the fair premium above the
current premium.  Specifically, when the premium crosses the accretion
threshold $\log(\mathrm{ELE})$ from below, the accumulation option
``switches on,'' creating a discrete increase in $\mathcal{F}$.

Combined with the concavity of $\mathcal{F}$ at high $\pi$ (diminishing
returns) and the flat region at low $\pi$ (no issuance), $G$ crosses zero
three times, yielding three fixed points.

\emph{Step 3: Stability.}
At $\pi_H^*$: $\mathcal{F}'(\pi_H^*) < 1$, so the equilibrium is stable
(perturbations are dampened by the mean-reversion of $\pi_t$).

At $\pi_L^*$: $\mathcal{F}'(\pi_L^*) < 1$ (the system is in the flat
region with no feedback), so this equilibrium is also stable.

At $\pi_U^*$: $\mathcal{F}'(\pi_U^*) > 1$ (the feedback amplifies
perturbations), making this equilibrium unstable.
\end{proof}

\subsubsection{Characterization of Equilibria}

\begin{proposition}[High-Premium Equilibrium]\label{prop:high_eq}
The high-premium equilibrium $\pi_H^*$ satisfies
\begin{equation}\label{eq:pi_H}
  \pi_H^* \approx \log\!\left(1 + \frac{g_B + \eta\, q_H / H}{r - g_B}
  \right),
\end{equation}
where $q_H = \alpha(\pi_H^* - \log(\mathrm{ELE}))^+$ is the equilibrium
accumulation rate.  The premium is sustained by:
\begin{enumerate}
  \item Accretive issuance (positive $g_B$), which grows BTC per share.
  \item Market impact ($\eta q_H / H$), which supports BTC price.
  \item Self-fulfilling expectations: the premium enables the issuance
    that justifies the premium.
\end{enumerate}
\end{proposition}

\begin{proof}
Substitute the market-impact-augmented BTC growth rate $\tilde{g}_B =
g_B + \eta q_H / H$ into the fair premium formula
\eqref{eq:pi_star_fp}.  The market impact term increases the effective
growth rate, which raises the present value of the accumulation option and
hence the premium.  The equilibrium is self-consistent because $q_H$ itself
depends on $\pi_H^*$.
\end{proof}

\begin{proposition}[Death Spiral Equilibrium]\label{prop:death_spiral}
The low-premium equilibrium $\pi_L^*$ is characterized by:
\begin{equation}\label{eq:pi_L}
  \pi_L^* \approx -\frac{\mathrm{PV}(\text{liability drag})}
  {\mathrm{NAV}_t^{\mathrm{ps}}},
\end{equation}
where the liability drag includes preferred dividend obligations $\Phi_t$
and coupon payments net of any residual operating cash flows.  The death
spiral mechanism operates as follows:
\begin{enumerate}
  \item Low premium $\pi_t < \log(\mathrm{ELE}_t)$ $\Rightarrow$
    no accretive ATM issuance possible.
  \item No accumulation $\Rightarrow$ no market impact support for $S_t$.
  \item BTC price stagnation or decline $\Rightarrow$ NAV falls.
  \item Ongoing liability payments (dividends, coupons) drain assets
    $\Rightarrow$ NAV falls further.
  \item Leverage $\ell_t$ rises $\Rightarrow$ accretion threshold
    $\log(1/(1-\ell_t))$ rises $\Rightarrow$ even harder to issue
    accretively.
  \item If the firm issues non-accretive preferred (last resort per
    Theorem~\ref{thm:optimal_channel}), $\Phi_t$ increases, adding
    more liability drag.
\end{enumerate}
\end{proposition}

\begin{proof}
When $\pi_t < \log(\mathrm{ELE}_t)$, the accumulation option value
$V_{\mathrm{acc}} = 0$ (the option is out-of-the-money with zero
probability of exercise under the stationary distribution).  The equity
value then equals NAV minus the present value of net liabilities:
\[
  E_t = \mathrm{NAV}_t - \mathrm{PV}(\text{liability drag}),
\]
so $\pi_L^* = \log(E_t / \mathrm{NAV}_t) < 0$ when the drag is positive.
Each step in the spiral raises $\ell_t$ or $\Phi_t$, increasing the drag
and deepening the discount, confirming local stability.
\end{proof}

\subsubsection{The Tipping Point}

\begin{definition}[Critical Premium]\label{def:critical_premium}
The \emph{critical premium} $\pi_c$ is the unstable equilibrium $\pi_U^*$
that separates the basins of attraction of the high-premium and death-spiral
equilibria:
\begin{equation}\label{eq:pi_critical}
  \pi_c = \pi_U^* = \inf\bigl\{\pi > \pi_L^* : G(\pi) = 0,\;
  G'(\pi) > 0 \bigr\}.
\end{equation}
\end{definition}

\begin{proposition}[Tipping Point Characterization]\label{prop:tipping}
The critical premium satisfies
\begin{equation}\label{eq:tipping_approx}
  \pi_c \approx \log(\mathrm{ELE}_t)
  + \frac{\sigma_\pi}{\sqrt{2\kappa}} \cdot
  \Phi_{\mathrm{N}}^{-1}\!\left(
    \frac{r - g_B}{g_B} \cdot \frac{e^{\pi_c} - 1}{1 + \sigma_\pi^2/(4\kappa r)}
  \right),
\end{equation}
and is approximately equal to the accretion threshold
$\log(\mathrm{ELE}_t)$ when market impact $\eta$ is small.
The tipping point has the following sensitivities:
\begin{align}
  \frac{\partial \pi_c}{\partial \bar\ell} &> 0
    &\quad& \text{(higher leverage raises the tipping point)},
    \label{eq:tip_ell} \\
  \frac{\partial \pi_c}{\partial \Phi_t} &> 0
    &\quad& \text{(higher dividend burden raises the tipping point)},
    \label{eq:tip_phi} \\
  \frac{\partial \pi_c}{\partial \sigma_S} &< 0
    &\quad& \text{(higher BTC vol lowers the tipping point)}.
    \label{eq:tip_sigma}
\end{align}
\end{proposition}

\begin{proof}
The tipping point is the premium at which the accumulation option value
just offsets the liability drag.  At $\pi_c$, the marginal value of a small
premium increase (which enables some accretive issuance) exactly equals the
marginal cost (the liability drag per unit time).  Setting
$V_{\mathrm{acc}}(\pi_c) = \mathrm{PV}(\text{drag})$ and solving for
$\pi_c$ using \eqref{eq:V_acc_closed} yields \eqref{eq:tipping_approx}.

The comparative statics follow from the implicit function theorem applied
to $G(\pi_c; \theta) = 0$:
\begin{itemize}
  \item Higher $\bar\ell$ raises $\log(\mathrm{ELE})$, requiring a higher
    premium to reach the accretive region.
  \item Higher $\Phi_t$ increases the liability drag that the accumulation
    option must offset.
  \item Higher $\sigma_S$ increases the option value for any given premium,
    lowering the premium needed to break even.  \qedhere
\end{itemize}
\end{proof}

\begin{remark}[Policy Implication]
The tipping point $\pi_c$ is the most important quantity for risk management.
Management should monitor the distance $\pi_t - \pi_c$: when this is large
and positive, the firm is safely in the high-premium basin.  When it
approaches zero, the firm is at risk of transitioning to the death spiral.
Issuing non-convertible preferred stock (which raises both $\bar\ell$ and
$\Phi_t$) pushes $\pi_c$ upward, shrinking the safety margin---a prediction
consistent with market reactions to STRF issuance.
\end{remark}

\subsubsection{Phase Diagram}

\begin{proposition}[Phase Diagram in $(\pi, S)$ Space]\label{prop:phase}
The dynamics of the endogenous premium system can be summarized in a
phase diagram with two regions separated by the critical surface
$\pi = \pi_c(\Phi_t)$:
\begin{enumerate}
  \item \textbf{Virtuous cycle region} ($\pi_t > \pi_c$):
    The drift of $\pi_t$ is toward $\pi_H^*$ from below, and the BTC price
    receives positive support from accumulation.  The expected trajectory is
    toward the high-premium equilibrium.
  \item \textbf{Death spiral region} ($\pi_t < \pi_c$):
    The drift of $\pi_t$ is toward $\pi_L^*$ from above, and BTC receives
    no accumulation support.  The expected trajectory is toward the discount
    equilibrium, with the approach accelerating as leverage rises.
\end{enumerate}
The velocity of the dynamics is:
\begin{equation}\label{eq:drift_pi}
  \mathbb{E}[d\pi_t \mid \pi_t = \pi] =
  \begin{cases}
    \kappa(\pi_H^* - \pi)\,dt & \text{if } \pi > \pi_c, \\[4pt]
    \kappa(\pi_L^* - \pi)\,dt & \text{if } \pi < \pi_c.
  \end{cases}
\end{equation}
\end{proposition}

\begin{proof}
This follows directly from the definition of $\pi_c$ as the boundary between
basins of attraction.  In each region, the premium mean-reverts toward the
locally stable equilibrium at speed $\kappa$.  At $\pi = \pi_c$ the drift
is zero (the unstable equilibrium), and any perturbation determines the
direction of subsequent evolution.
\end{proof}

%% ====================================================================
\subsection{Mispricing Dynamics}\label{subsec:mispricing}

We now define the deviation of the observed premium from its fair value and
characterize its dynamics.

\begin{definition}[Mispricing]\label{def:mispricing}
The \emph{mispricing} at time $t$ is the difference between the observed
log-premium and the fair log-premium:
\begin{equation}\label{eq:delta_def}
  \Delta_t = \pi_t - \pi^*(\Phi_t).
\end{equation}
Positive $\Delta_t$ indicates overvaluation (the market assigns a higher
premium than fundamentals justify).  Negative $\Delta_t$ indicates
undervaluation.
\end{definition}

\begin{theorem}[Mispricing Dynamics]\label{thm:mispricing_dynamics}
Under the endogenous premium system
\eqref{eq:dS_endo}--\eqref{eq:dH_endo}, the mispricing $\Delta_t$ follows
\begin{equation}\label{eq:d_delta}
  d\Delta_t = -\kappa_\Delta\, \Delta_t\, dt
  + \sigma_\Delta\, d\tilde{W}_t,
\end{equation}
where
\begin{align}
  \kappa_\Delta &= \kappa + \frac{\partial \pi^*}{\partial \Phi}
    \cdot \frac{\partial \Phi}{\partial \pi}\bigg|_{\pi = \pi^*}
    > \kappa,
    \label{eq:kappa_delta} \\[4pt]
  \sigma_\Delta^2 &= \sigma_\pi^2
    + \left(\frac{\partial \pi^*}{\partial S}\right)^2 \sigma_S^2 S_t^2
    - 2\rho\, \sigma_\pi \cdot \frac{\partial \pi^*}{\partial S}\,
    \sigma_S S_t,
    \label{eq:sigma_delta}
\end{align}
and $\tilde{W}_t$ is a Brownian motion under $\mathbb{P}$.  The key result
is that $\kappa_\Delta > \kappa$: \emph{mispricing mean-reverts faster than
the premium itself}.
\end{theorem}

\begin{proof}
Apply It\^{o}'s lemma to $\Delta_t = \pi_t - \pi^*(\Phi_t)$:
\begin{equation}\label{eq:ito_delta}
  d\Delta_t = d\pi_t - d\pi^*(\Phi_t).
\end{equation}
The first term is given by \eqref{eq:dpi_endo}:
\[
  d\pi_t = \kappa(\pi^* - \pi_t)\,dt + \sigma_\pi\,dW_t^\pi
  = -\kappa\,\Delta_t\,dt + \sigma_\pi\,dW_t^\pi.
\]
For the second term, $\pi^*$ depends on $\Phi_t$ which evolves with
$S_t$ and $H_t$.  By the chain rule:
\[
  d\pi^* = \frac{\partial \pi^*}{\partial S}\,dS_t
  + \frac{\partial \pi^*}{\partial H}\,dH_t + \cdots
\]
The crucial feedback term is $(\partial\pi^*/\partial\Phi) \cdot
(\partial\Phi/\partial\pi)$: when the premium is above $\pi^*$, the
resulting accumulation updates fundamentals, which moves $\pi^*$
\emph{toward} the current premium.  This reduces the mispricing from both
sides---$\pi_t$ mean-reverts down (by OU dynamics) while $\pi^*$ adjusts up
(by fundamental improvement)---yielding the enhanced mean-reversion speed
$\kappa_\Delta > \kappa$.

Collecting the drift terms:
\[
  \mathbb{E}[d\Delta_t \mid \Delta_t]
  = -\kappa\,\Delta_t - \frac{\partial\pi^*}{\partial\Phi}
  \cdot \frac{\partial\Phi}{\partial\pi}\,\Delta_t
  = -\kappa_\Delta\,\Delta_t,
\]
where the second term is positive (improvement in fundamentals when $\pi >
\pi^*$), confirming \eqref{eq:kappa_delta}.  The volatility
\eqref{eq:sigma_delta} follows from computing $\mathrm{Var}(d\Delta_t)$
using the correlated Brownian motions.
\end{proof}

\begin{proposition}[Stationary Distribution of Mispricing]%
\label{prop:mispricing_stat}
Under the linearized dynamics \eqref{eq:d_delta}, the mispricing has a
stationary distribution
\begin{equation}\label{eq:delta_stationary}
  \Delta_\infty \sim \mathcal{N}\!\left(0,\;
  \frac{\sigma_\Delta^2}{2\kappa_\Delta}\right).
\end{equation}
The stationary variance of mispricing is smaller than the stationary
variance of the premium:
\begin{equation}\label{eq:var_comparison}
  \frac{\sigma_\Delta^2}{2\kappa_\Delta}
  < \frac{\sigma_\pi^2}{2\kappa},
\end{equation}
whenever the fundamental feedback term
$(\partial\pi^*/\partial\Phi)(\partial\Phi/\partial\pi) > 0$.
\end{proposition}

\begin{proof}
The OU process \eqref{eq:d_delta} with mean zero, speed $\kappa_\Delta$,
and volatility $\sigma_\Delta$ has stationary variance
$\sigma_\Delta^2/(2\kappa_\Delta)$.  Since $\kappa_\Delta > \kappa$ and
$\sigma_\Delta \le \sigma_\pi + |\partial\pi^*/\partial S|\sigma_S S_t$
(with the cross term providing reduction when $\rho$ has the appropriate
sign), the variance ratio satisfies
\[
  \frac{\sigma_\Delta^2 / (2\kappa_\Delta)}{\sigma_\pi^2 / (2\kappa)}
  = \frac{\sigma_\Delta^2}{\sigma_\pi^2} \cdot
  \frac{\kappa}{\kappa_\Delta}
  < 1
\]
when the feedback effect dominates the additional volatility from the
moving $\pi^*$.  The condition
$(\partial\pi^*/\partial\Phi)(\partial\Phi/\partial\pi) > 0$ ensures that
the enhanced mean-reversion effect is large enough.
\end{proof}

\begin{remark}[Trading Implications]\label{rem:trading}
The result $\kappa_\Delta > \kappa$ has a concrete trading implication:
a mean-reversion strategy based on mispricing $\Delta_t$ (trading toward
the model-implied fair premium) should converge faster than a strategy
based on the raw premium $\pi_t$ (trading toward the historical mean
$\theta$).  The mispricing signal $\Delta_t$ is informationally superior
because it incorporates the current state of fundamentals, while the raw
premium signal compares against a static target.
\end{remark}

\begin{definition}[Mispricing Z-Score]\label{def:mispricing_zscore}
Analogous to the PMRI of Section~\ref{sec:methodology}, we define the
\emph{Mispricing Z-Score}:
\begin{equation}\label{eq:mzs}
  \mathrm{MZS}_t = \frac{\Delta_t}{\sigma_\Delta / \sqrt{2\kappa_\Delta}}
  = \frac{\pi_t - \pi^*(\Phi_t)}{\sigma_\Delta / \sqrt{2\kappa_\Delta}}.
\end{equation}
Under the stationary distribution, $\mathrm{MZS}_t \sim \mathcal{N}(0,1)$.
Values $|\mathrm{MZS}_t| > 2$ indicate significant mispricing
(approximately 5\% tail probability).
\end{definition}

\begin{remark}[Relationship to PMRI]
The PMRI standardizes $\pi_t - \theta$ against $\sigma_\pi/\sqrt{2\kappa}$,
using the exogenous long-run mean.  The MZS standardizes $\pi_t - \pi^*$
against $\sigma_\Delta/\sqrt{2\kappa_\Delta}$, using the endogenous fair
premium.  The MZS is a strictly more informative signal: it can distinguish
between (i) a high premium that is justified by strong fundamentals
($\mathrm{MZS} \approx 0$, $\mathrm{PMRI} \gg 0$) and (ii) a high premium
that reflects bubble dynamics ($\mathrm{MZS} \gg 0$, $\mathrm{PMRI} \gg 0$).
The PMRI conflates both cases.
\end{remark}

%% ====================================================================
\subsection{Summary of Key Results}\label{subsec:endo_summary}

\begin{table}[htbp]
\centering
\caption{Summary of endogenous premium theory results.}
\label{tab:endo_summary}
\small
\begin{tabular}{@{}lp{8cm}l@{}}
\toprule
\textbf{Result} & \textbf{Statement} & \textbf{Ref.} \\
\midrule
Fair premium & $\pi^* = \log(1 + V_{\mathrm{acc}} / \mathrm{NAV}^{\mathrm{ps}})$;
  unique fixed point & Thm.~\ref{thm:fair_premium} \\
Volatility is value & $\partial\pi^*/\partial\sigma_S > 0$: higher BTC vol
  increases fair premium & Prop.~\ref{prop:comp_statics} \\
Multiple equilibria & High-premium (stable), death spiral (stable),
  tipping point (unstable) & Thm.~\ref{thm:multiple_eq} \\
Tipping point & $\pi_c \approx \log(\mathrm{ELE})$; rises with leverage
  and dividend burden & Prop.~\ref{prop:tipping} \\
Mispricing reverts faster & $\kappa_\Delta > \kappa$: fundamental feedback
  accelerates convergence & Thm.~\ref{thm:mispricing_dynamics} \\
MZS signal & Strictly more informative than PMRI for detecting over/under-valuation
  & Def.~\ref{def:mispricing_zscore} \\
\bottomrule
\end{tabular}
\end{table}

\noindent The endogenous premium framework transforms the NAV premium from
an unexplained market phenomenon into a derived quantity with clear
fundamental drivers.  It provides testable predictions (volatility--premium
correlation, mispricing mean-reversion speed, tipping point location) and
a superior valuation signal (MZS vs.\ PMRI) for practitioners.
